%% T1.tex -- Proposition 1: Closed-form reliability via the fundamental matrix.

\noindent\emph{The following proposition is a direct adaptation of the
Kemeny--Snell fundamental-matrix identity \cite[Ch.~3]{kemenysnell} for
absorbing DTMCs, specialized to the agent state space
(Definition~\ref{def:amc}). No new mathematical content is claimed.}

\begin{proposition}[Closed-Form Reliability; adapted from Kemeny--Snell \cite{kemenysnell}]
\label{thm:closed}
Under Assumptions~\ref{A:trans}--\ref{A:init}, for every $d\in\mathbb{N}_{\ge 0}$,
\begin{equation}
R(d) \;=\; e_{s_0}^\top(I-Q^{d})\,N\,R_\oplus,
\qquad N \triangleq (I-Q)^{-1},
\label{eq:t1finite}
\end{equation}
and
\begin{equation}
R_\infty \;=\; \lim_{d\to\infty}R(d) \;=\; e_{s_0}^\top\,N\,R_\oplus.
\label{eq:t1inf}
\end{equation}
\end{proposition}

%% ----------------------------------------------------------------
%% [REDUCTION 2026-04-24] Proof and remarks retained in source but
%% suppressed from the 12-page body. Full proof is in the
%% reproducibility artifact at proofs/T1.tex (this file).
%% ----------------------------------------------------------------
\iffalse
\begin{proof}
\emph{Step 1 — telescoping.}
For any square matrix $Q$ and $d\ge 1$,
\[
(I-Q)\sum_{t=0}^{d-1}Q^{t}
  \;=\;\sum_{t=0}^{d-1}Q^{t}-\sum_{t=1}^{d}Q^{t}
  \;=\;I-Q^{d}.
\]
\emph{Step 2 — inversion.}
By A\ref{A:inv}, $(I-Q)$ is invertible, so
\[
\sum_{t=0}^{d-1}Q^{t} \;=\; (I-Q)^{-1}(I-Q^{d}) \;=\; N(I-Q^{d}).
\]
Substituting into Definition~\ref{def:reliab} gives
$R(d)=e_{s_0}^\top N(I-Q^{d})R_\oplus$. Since $N$ and $(I-Q^{d})$ are $m\times m$
matrices that commute
(because $N$ is a polynomial in $Q$), the product may be written as
$(I-Q^{d})N$, establishing \eqref{eq:t1finite}.

\emph{Step 3 — limit.}
A\ref{A:trans} gives $\rho(Q)<1$, so $\|Q^{d}\|\to 0$ as $d\to\infty$ in any
sub-multiplicative norm. Therefore $(I-Q^{d})\to I$ entry-wise, and
$R(d)\to e_{s_0}^\top N R_\oplus$, which is \eqref{eq:t1inf}.
\end{proof}

\begin{remark}
The entry $N_{ij}$ has the classical interpretation (Kemeny--Snell) as the
expected number of visits to transient state $j$ starting from state $i$ before
absorption. Equation~\eqref{eq:t1inf} is therefore a sum over transient
states of ``visits $\times$ per-step success probability'', which provides
a sanity check of the algebra.
\end{remark}

\begin{remark}[Complexity]
Computing $R_\infty$ by \eqref{eq:t1inf} costs $O(m^3)$ for the inversion,
dominated by the LU decomposition of $(I-Q)$. Computing $R(d)$ for a
schedule of horizons $d\in\{d_1,\ldots,d_L\}$ via repeated squaring of $Q$
costs $O(m^3\log d_{\max})$.
\end{remark}
\fi
